\section{The periodic hitting problem}\label{sec:periodic}

For a measurable $1$-periodic set $A\subseteq\R$, write
\[
 \rho(A)=m(A\cap[0,1])
\]
for its density in one period. We will sometimes identify such a set
with its image in the circle $\T=\R/\mathbb Z$, equipped with normalized
Lebesgue measure. This identification preserves the density. Open or
closed subsets of the circle lift to open or closed periodic subsets of
$\R$.

\begin{lemma}[Periodic hitting sets]\label{lem:periodic}
For every $p\in(0,1)$, there exists an open $1$-periodic set $H\subseteq\R$
such that $\rho(H)\le6p$ and
\[
 \forall x\in\R\ \forall t\in[1,2]\ \exists n\in\mathbb Z_{\ge1}:
 \qquad x+t2^{-n}\in H.
\]
\end{lemma}

The same set $H$ works for every center and every scale in the displayed
interval. The witness $n$ may depend on both. In
Section~\ref{sec:global}, we apply the lemma with a summable sequence of
parameters, rescale the resulting periodic sets by powers of two, and
include their reflections. This handles every nonzero real scale while
keeping the total measure removed from $[0,1]$ small.

For the proof of Lemma~\ref{lem:periodic}, fix $p\in(0,1)$. We construct
a random periodic set $B$ with expected density $p$, together with a
finite set $\mathcal N$ of positive integer indices. Outside a
deterministic set of centers of density less than $p$, the probability
that some $t\in[1,2]$ avoids $B$ at all indices in $\mathcal N$ will be
at most $2p$. We then enlarge $B$ to an open set, bound the expected
measure of its exceptional centers, and cover those centers by a small
open set. The last step uses further indices in the dyadic tail;
$\mathcal N$ need not provide the hits at the repaired centers.
